\section{The Exercise Indicator (Supplementary)}
\label{app:exercise_indicator}

\subsection{Scope limitation}

The exercise indicator is retained as a \emph{supplementary}
measurement tool, applicable only where interior access is
legitimate---specifically, within the
collective's own operational perimeter where no privacy issue
arises (e.g., the framework measuring the exercise status of its
own internal components, or a cooperative measurement within a
consenting sub-collective).

\begin{definition}[Exercise Indicator, Internal Perimeter]
\label{def:exercise_indicator}
Let $G$ be a capability collective operating within an
agreed-upon observation perimeter $O \subseteq P$.  We require
$O$ to be \emph{boundary-closed in the poset-independence sense
of Definition~\ref{def:category_partition}(i)}: neither
cooperative-composition nor subsumption links cross the
$O$/$P \setminus O$ partition.  Concretely, two conditions:
\begin{enumerate}
\item[(BC1)] \emph{Cooperative closure.}  For every cooperative
  capability $z \in P$, either all participants of $z$ lie in
  $O$ or none do (so cooperatives are not split across the
  boundary).
\item[(BC2)] \emph{Subsumption closure.}  No subsumption
  relation crosses the boundary: if $c \succeq c'$ in $P$ (i.e.,
  $c$ subsumes $c'$), then either both $c, c' \in O$ or both
  $c, c' \in P \setminus O$.
\end{enumerate}
Together (BC1) and (BC2) make $O$ and $P \setminus O$
poset-independent in the sense of
Definition~\ref{def:category_partition}(i), which is what
licenses additive composition across the boundary below.  (BC1)
alone---the weaker cooperative-only condition---is insufficient:
subsumption edges crossing the boundary would break M4 on the
$O$/$P \setminus O$ partition.

For each capability $d \in O$, define:
\[
  e_t(d) = \begin{cases}
  1 & \text{if there exists a realized cooperative output in
    $[t - \Delta, t]$} \\
    & \text{such that removing $d$ from $O$ makes the output
      unrealizable,} \\
  0 & \text{otherwise.}
  \end{cases}
\]
The indicator is scoped to the observation perimeter $O$.  The
framework may use $e_t(d)$ for internal diagnostics within $O$,
but may \emph{not} apply it to capabilities outside $O$ (where
privacy constraints apply).

The \emph{exercise-based realized volume} within $O$ is the
$\volL$-weighted sum of exercised capabilities in the
perimeter:
\[
  \volRexact(O) \;:=\;
  \volL\bigl(\{d \in O : e_t(d) = 1\}\bigr).
\]
For the window-integrated form used by
Theorem~\ref{thm:aggregate_bound}, we set
\[
  \volRexact^{[W]}(O) \;:=\;
  \volL\bigl(O \cap P^{\mathrm{ex}}_W\bigr)
\]
where $P^{\mathrm{ex}}_W = \bigcup_{t \in W} \{d \in O : e_t(d)
= 1\}$.  The instantaneous form $\volRexact(O)$ is the
$t$-indexed specialization; references to $\volRexact$ without a
window subscript refer to the window-integrated form under the
same convention as Definition~\ref{def:volR}.
\end{definition}

\begin{remark}[The two channels are complementary]
Within the perimeter $O$, the exercise indicator provides exact
exercise status ($e_t(d) \in \{0, 1\}$).  Outside $O$, the
sacrifice channel provides a lower bound.  Under the
boundary-closed hypothesis of
Definition~\ref{def:exercise_indicator}, the aggregate $\volR$
diagnostic is a valid lower bound obtained by additive
composition:
\[
  \volRcomb \;:=\; \volRexact(O) + \volRlower(P \setminus O)
  \;\leq\; \volR(O) + \volR(P \setminus O)
  \;=\; \volR(P),
\]
where the final equality uses axiom M4 applied to the
poset-independent pair $O$, $P \setminus O$
(Section~\ref{sec:axiom_inheritance}).  $\volRexact(O)$ uses the
exercise indicator and $\volRlower(P \setminus O)$ uses the
sacrifice channel.
\end{remark}

\begin{remark}[When boundary-closure fails]
\label{rem:boundary_cooperative_undercount}
If $O$ fails either condition---a cooperative $z$ has
participants split across $O$ and $P \setminus O$ (BC1 fails),
or a subsumption edge $c \succeq c'$ crosses the boundary (BC2
fails)---then $O$ and $P \setminus O$ are not
poset-independent, and M4 additivity across the boundary does
not apply.  Specifically: under a BC1 failure, $z$ contributes
to $\volR(P)$ but to neither $\volRexact(O)$ nor $\volRlower(P
\setminus O)$ individually (the exercise indicator inside $O$
evaluates only capabilities fully within $O$; the sacrifice
channel outside $O$ cannot reach $z$ if $z$'s full participant
set is not in $P \setminus O$).  Under a BC2 failure, the
subsumption edge couples $\volL$-mass across the boundary in
ways that the independent-component additivity of M4 cannot
decompose.  In either case $\volRcomb$ remains a
\emph{conservative} lower bound on $\volR(P)$---the subadditivity
of $\volL$ on non-independent components ensures the sum of
pieces is $\leq$ the whole---but the equality with $\volR(P)$
above no longer holds.  Practical perimeters (the collective's
own operational scope, consenting sub-collectives) typically
satisfy both boundary-closure conditions because cooperative
capabilities and subsumption hierarchies are either wholly
internal to the perimeter or wholly external; where either BC1
or BC2 fails, the undercount is quantified by the $\volL$-weight
of the offending cross-boundary cooperatives and subsumption
edges.
\end{remark}
